% SPDX-License-Identifier: CC-BY-SA-4.0
% Part: Introduction
% Section: Words and languages
% Exercise: Languages and operators

\begingroup

\begin{exercice}[Langages et opérateurs]\label{exo:introduction/languages/languages}

  Soit l'alphabet $\Sigma = \{a, b\}$ et les langages :
  \begin{itemize}
  \item $L_1 = \{a, ab, ba\}$,
  \item $L_2 = \{\varepsilon, b, ab\}$,
  \item $L_3 = \{u \in \Sigma^\star \mid |u|_a = |u|_b \}$,
  \item $L_4 = \{u \in \Sigma^\star \mid |u| \mod 2 = 1\}$,
  \item $L_5 = \{u \in \Sigma^\star \mid \exists n \in \mathbb{N}, \exists m\in \mathbb{N}, u = a^n \cdot b^m \} = \{a^n \cdot b^m \mid n \ge 0, m \ge 0\}$.
  \end{itemize}

  \begin{question}
  \item Donner, en français, la définition des langages $L_3$, $L_4$ et $L_5$.
  \item Donner les résultats des opérations suivantes :
    $$\begin{array}{c@{\quad\quad}c@{\quad\quad}c@{\quad\quad}c@{\quad\quad}c@{\quad\quad}c}
      L_1 \cup L_2                & 
      L_1 \cap L_2                & 
      L_3 \cap L_4                & 
      L_3 \cap L_5                & 
      \Sigma^\star \setminus L_5   & 
      L_1 \cdot L_2               \\[1mm]
      L_2 \cdot L_1               & 
      L_1^2                       & 
      L_2^3                       & 
      L_3^\star                    &
      L_1 \cup \emptyset          & 
      L_2 \cup \{\varepsilon\}    \\[1mm] 
      L_1 \cap \emptyset          & 
      L_2 \cap \{\varepsilon\}    & 
      L_1 \cdot \emptyset         & 
      \emptyset \cdot L_1         &
      L_1 \cdot\{\varepsilon\}    &
      \{\varepsilon\}\cdot L_1     
    \end{array}$$
  \end{question}
  Soit $L$ et $M$ deux langages. Que peut-on dire de $L$ et $M$ dans les cas suivants : 
  \begin{question}
  \item si $L \cdot M = \{\varepsilon\}$ ?
  \item si $L \cdot M = \emptyset$ ?
  \end{question}

\end{exercice}

\begin{correction}

  \begin{question}
  \item
    \begin{itemize}
    \item $L_3$ est l'ensemble des mots qui contiennent autant de $a$ que de $b$.
    \item $L_4$ est l'ensemble des mots de longueur impaire.
    \item $L_5$ est l'ensemble de mots constitués d'une suite de $a$, éventuellement
      vide, suivie d'une suite de $b$, éventuellement vide.
    \end{itemize}
  \item
    $\begin{array}[t]{rll}
    L_1 \cup L_2                &=& \{\varepsilon, a, b, ab, ba\}   \\ 
    L_1 \cap L_2                &=& \{ab\}   \\ 
    L_3 \cap L_4                &=& \emptyset   \\ 
    L_3 \cap L_5                &=& \{a^n b^n \mid n\in \mathbb{N}\}   \\ 
    \Sigma^\star \setminus L_5   &=& \{u \cdot ba \cdot v | u , v \in \Sigma^\star\}   \\ 
    L_1 \cdot L_2               &=& \{a, ab, ba, abb, bab, aab, abab, baab\}   \\ 
    L_2 \cdot L_1               &=& \{a, ba, aba, ab, bab, abab, bba, abba\}  \\ 
    L_1^2                       &=& \{aa, aab, aba, abab, abba, baa, baab, baba\}\\
    L_2^3                       &=& \{\varepsilon, b, ab, bb, bab, abb, abab\} \cdot \{\varepsilon, b, ab\}\\
    &=& \{\varepsilon, b, ab, bb, bab, abb, abab, bbb, babb, abbb, ababb, bbab, babab, abbab, ababab\}\\
  \end{array}$

    $$\begin{array}[t]{rll@{\quad\quad\quad}rll@{\quad\quad\quad}rll}
    L_3^\star                    &=& L_3             &
    L_1 \cup \emptyset          &=& L_1             & 
    L_1 \cup \{\varepsilon\}    &=& L_1             \\ 
    L_1 \cap \emptyset          &=& \emptyset       & 
    L_1 \cap \{\varepsilon\}    &=& \{\varepsilon\} & 
    L_1 \cdot \emptyset         &=& \emptyset       \\ 
    \emptyset \cdot L_1         &=& \emptyset       & 
    L_1 \cdot\{\varepsilon\}    &=& L_1             & 
    \{\varepsilon\}\cdot L_1    &=& L_1             \\ 
  \end{array}$$
  \item $L = \{\varepsilon\}$ et $M = \{\varepsilon\}$.
  \item $L = \emptyset$ ou $M = \emptyset$.
  \end{question}

\end{correction}

\endgroup
\endinput
